% 原模板的空白正文页。以下命令只登记目录锚点，不打印标题。
% 可在相应命令后填写正文，目录页码自动更新；默认仍为空白页。
\TemplateSection{1.\quad 启发式算法综述}
\TemplateSubsection{1.1 算法描述}
% 在此填写。
\TemplateSubsection{1.2 算法的可采纳性}
\TemplateSubsubsection{1.2.1 A*算法}
\TemplateSubsubsection{1.2.2 A*算法可采纳性证明}
\TemplateSubsection{1.3 重复扩展结点的问题}
\TemplateSubsubsection{1.3.1 重复扩展结点问题}
\TemplateSubsubsection{1.3.2 启发函数的单调性}
\TemplateSubsection{1.4 IDA*算法基本原理及其可采纳性}

\TemplateSection{2.N-Puzzle游戏}
\TemplateSubsection{2.1 N-Puzzle游戏描述}
\TemplateSubsection{2.2 游戏的建模，状态及状态空间，状态转换规则}
\TemplateSubsection{2.3 N-Puzzle可解性研究}
\TemplateSubsection{2.4 三个启发函数的有效性和单调性证明}
\TemplateSubsubsection{2.4.1 曼哈顿距离}
\TemplateSubsubsection{2.4.2 不在位将牌}
\TemplateSubsubsection{2.4.3 Disjoint Pattern}

\TemplateSection{3. N-Puzzle求解算法实现}
\TemplateSubsection{3.1 状态及结点的数据结构}
\TemplateSubsection{3.2 Open表及Closed表的数据结构}
\TemplateSubsection{3.3 三个启发函数的具体实现}
\TemplateSubsubsection{3.3.1 曼哈顿距离}
\TemplateSubsubsection{3.3.2 不在位将牌}
\TemplateSubsubsection{3.3.3 Disjoint Pattern和Pattern Database}
\TemplateSubsection{3.4 求解后继结点的函数}

\TemplateSection{4.N-Puzzle问题求解探讨{\color{red}（提供详细的实验数据分析）}}
\TemplateSubsection[templatetwoextra]{4.1 Open表及Closed表数据结构的改进及其对算法性能的影响}
\TemplateSubsection{4.2 A*算法，IDA*算法的性能比较}
\TemplateSubsection{4.3 三个不同的启发函数及IDA*的性能比较}
\TemplateSubsubsection{4.3.1 运行时间对比及分析}
\TemplateSubsubsection{4.3.2 扩展结点比较及分析}

\TemplateSection{5.解路径的可视化}

% 原目录的“总结及展望”前只有一个点，照原样保留，不擅自补编号。
\TemplateSection{.总结及展望}

\TemplateSection{参考文献}
